\chapter{Reproduction Recipe}
\label{AppendixC}
\lhead{Appendix C. \emph{Reproduction Recipe}}

Every quantitative result in Chapter~\ref{Chapter6} is regenerated by
the command sequence below on a clean Ubuntu~24.04 host running Linux
kernel~6.17 and BCC~0.31. The full run is driven by the project
\texttt{Makefile}; each target below maps to one step.

\begin{verbatim}
# 1. Prerequisites (one-time)
$ sudo apt-get install -y build-essential bpfcc-tools \
      python3-bpfcc iproute2 docker.io docker-compose-v2 wrk falco
$ python3 -m venv .venv && source .venv/bin/activate
$ pip install -r requirements.txt

# 2. Bring the 22-container testbed up
$ make up                                 # docker compose up + routes

# 3. Calibrate (writes calibration/*.npz, *.pkl, *.json)
$ make calibrate

# 4. Full marathon (CausalTrace + Falco stock/tuned + Tetragon stock/tuned,
#    one seed-42 permutation of 155 attacks replayed per detector)
$ make marathon

# 5. Analyse and regenerate every paper figure + Wilson-CI table
$ python3 scripts/marathon_analyze.py results/marathon
$ python3 generate_astar_plots.py --data-dir results/marathon \
      --out-dir BTech_Project_Report/figures
$ python3 BTech_Project_Report/figures/build_detection_table.py

# 6. Rebuild the thesis PDF
$ cd BTech_Project_Report && pdflatex main && bibtex main && \
      pdflatex main && pdflatex main
\end{verbatim}

The seeded attack permutation is stored in
\texttt{results/marathon/state.json}; resumable runs use
\texttt{run\_marathon\_evaluation.py --resume}. The detection-window
width ($\Delta = 25$\,s) and the CCA rank ($k = 15$) are the only
hyper-parameters that need to be kept consistent between calibration and
evaluation; everything else is data-driven.
